% Exact positivity reductions; owned by cycles. No preamble.
\paragraph{C.1o (Proved: projective-space stabilization identifies the exact positivity kernel).}
Let $X/\mathbb C$ be smooth projective of dimension $n$, let $h$ be an
ample rational divisor class, and put $t=c_1(\mathcal O_{\mathbb P^d}(1))$.
Suppose that $n+d=2m$ is even. For $0\le d\le n$ set
$i=(n-d)/2$. There is an explicit isomorphism
\begin{equation}\label{eq:C1o-radical}
 \begin{split}
 \ker\bigl(h^{d+1}:G^i(X)\longrightarrow G^{i+d+1}(X)\bigr)
 &\xrightarrow{\ \sim\ }
 \operatorname{rad}\bigl(q_{X\times\mathbb P^d,h+t}\bigr),\\
 z&\longmapsto
 \sum_{j=0}^{d}(-1)^j h^jz\,t^{d-j}.
 \end{split}
\end{equation}
Here $q$ is the middle primitive form of C.1j, with cycles modulo
algebraic equivalence and rational coefficients. If $d>n$, that radical
is zero. The statement uses the actual product polarization $h+t$;
it is unchanged after a positive rational rescaling of either factor.

\begin{proof}
The projective-bundle formula descends to both algebraic and
homological equivalence. It gives
\[
 G^m(X\times\mathbb P^d)=
       \bigoplus_{b=0}^{d}G^{m-b}(X)t^b.
\]
Write $\Gamma=\sum_{b=0}^d a_bt^b$ in this direct sum.
The equation $(h+t)\Gamma=0$ is exactly
\[
 ha_0=0,\qquad ha_b+a_{b-1}=0\quad(1\le b\le d).
\]
Hence all coefficients are determined by $z=a_d$:
$a_{d-j}=(-h)^jz$, and the remaining equation is $h^{d+1}z=0$.
Its degree is $m-d=(n-d)/2=i$. Conversely those formulas give
\[
 (h+t)\sum_{j=0}^{d}(-1)^j h^jz\,t^{d-j}
                         =(-1)^d h^{d+1}z,
\]
because $t^{d+1}=0$. The coefficient of $t^d$ recovers $z$, so
the map is injective as well as surjective. If $d>n$, its putative
last coefficient has negative codimension and vanishes; the same
recursion then kills every coefficient. C.1j identifies this primitive
Griffiths kernel with the radical, proving the claim. For a polarization
$h+bt$ with $b>0$, divide by $b$ and apply the same argument to
$h/b+t$. Positive rescaling does not change the kernel or strict
positivity.
\end{proof}

It follows that the following two \emph{universal} assertions over
$\mathbb C$ are equivalent:
\begin{enumerate}
\item the letter's strict middle positivity $B$ holds for every smooth
projective variety of even dimension and every ample class;
\item for every smooth projective $n$-fold $X$, every ample $h$, and
every $0\le 2i\le n$, the map
\begin{equation}\label{eq:C1o-shifted}
 h^{n-2i+1}:G^i(X)\longrightarrow G^{n-i+1}(X)
 \quad\text{is injective}.
\end{equation}
\end{enumerate}
For the forward implication take $d=n-2i$ in
\eqref{eq:C1o-radical}; for the reverse implication take $d=0$.
The forward implication also gives the injectivity part of the
letter's $A$: if $h^{n-2i}x=0$ modulo algebraic equivalence,
cohomological hard Lefschetz first puts $x$ in $G^i(X)$, and
\eqref{eq:C1o-shifted} then kills it. No surjectivity assertion follows
from this argument.

There is an unconditional and useful test family. For every smooth
projective threefold $V$, every odd $d\ge1$, and every ample class on
$V\times\mathbb P^d$, the letter's $B$ holds. Indeed, the only
nonnegative degrees $i=(3-d)/2$ in the formula are $i=1$ and $i=0$,
and $G^1(V)=G^0(V)=0$. Any ample class on this product has the form
$h+bt$ with $h$ ample on $V$ and $b>0$, by the projective-bundle
description of its Picard group and restriction to its two factors.
Taking the specified quintic threefold of C.1g therefore gives, in
every even dimension at least four, a variety satisfying strict
middle positivity while possessing a nonzero Griffiths group and
a nonzero radical in its correspondence algebra by C.1l.
With a very general quintic, Clemens's theorem in C.1k strengthens
the Griffiths group to infinite dimension. This refutes the
\emph{individual-variety} inference from its own middle positivity
to the corresponding conclusions. It does not refute
an implication whose premise is positivity for \emph{all} varieties
and products; that universal premise remains precisely the
unproved assertion \eqref{eq:C1o-shifted}.

\paragraph{C.1p (Proved: the exact Beauville-graded positivity test on abelian stabilizations).}
Fix a complex abelian variety $A$ of dimension $g$. The following
assertions are equivalent:
\begin{enumerate}
\item the letter's $B$ holds for every even-dimensional
$A\times\mathbb P^d$ and every ample class on it;
\item $G^r(A)_{(s)}=0$ for every $r$ and every $s\le0$.
\end{enumerate}
Equivalently, the quotient
$B^*(A)=\mathrm{CH}^*(A)_{\mathbb Q}/
\mathrm{CH}^*(A)_{\mathrm{alg},\mathbb Q}$ has no negative
Beauville pieces, and its grade-zero cycle map is injective.
This is a statement modulo algebraic equivalence. We do not identify
it with the corresponding conjectures for the ordinary Chow group.

\begin{proof}
We use Beauville's algebraic $\mathfrak{sl}_2$-action, as in C.3h
\citep[Theorem 4.2 and Proposition 5.2, pp.~260--261]{rootBeauville2010}.
His footnote 1 on p.~254 permits any adequate equivalence.
The correspondences preserve the homologically trivial subgroup and
algebraic equivalence, so $G^*(A)$ is a direct sum of finite
dimensional strings
\[
 z,\theta z,\ldots,\theta^{g+s-2a}z,
 \qquad z\in G^a(A)_{(s)},\quad Fz=0,
 \quad g+s-2a\ge0,
\]
where $\theta$ is the symmetric ample divisor and $F$ is the lowering
operator. If a nonzero piece with $s\le0$ exists, choose a nonzero
lowest vector in one of its strings. Then $2a\le g$, and with
$d=g-2a$ we have
\[
                   \theta^{d+1}z=0,
       \qquad g+s-2a\le d.
\]
Formula \eqref{eq:C1o-radical} supplies a nonzero primitive radical
on $A\times\mathbb P^d$, disproving (1).

Conversely suppose (2). Every remaining string has $s\ge1$.
For a degree-$r$ element $\theta^{r-a}z$, the power
$\theta^{g-2r+1}$ could kill it only if
\[
 r-a+g-2r+1>g+s-2a,
 \quad\text{that is,}\quad a>r+s-1.
\]
But $a\le r$ and $s\ge1$, so this is impossible. The direct sum of
strings therefore gives all injectivities
\eqref{eq:C1o-shifted} for this $A$, and C.1o proves (1).
An ample divisor and its symmetric component have the same class
modulo algebraic equivalence, so this also handles nonsymmetric
ample classes. The description of ample classes on a product with
projective space handles all the polarizations in (1).

For $s\ne0$, a Beauville eigenclass has zero cohomology class:
multiplication by an integer $q>1$ acts on its cycle class both by
$q^{2r-s}$ and by $q^{2r}$. Thus $G^r(A)_{(s)}=B^r(A)_{(s)}$
for $s\ne0$, while $G^r(A)_{(0)}$ is exactly the kernel of the
grade-zero cycle map. This proves the equivalent formulation.
\end{proof}

The same string argument gives precise first tests without assuming
any Beauville conjecture. For $g\le3$, every string with $s\le0$
would have lowest degree $a\le1$, where the Griffiths group is zero;
hence (1) holds unconditionally. For an abelian fourfold,
\[
 \operatorname{rad}(q_{A,\theta})=G^2(A)_{(0)}.
\]
For an abelian fivefold,
\[
 \operatorname{rad}(q_{A\times\mathbb P^1,\theta+t})
       \simeq G^2(A)_{(-1)}\oplus G^2(A)_{(0)}.
\]
Indeed the only possible lowest degree is $a\ge2$: in dimension
four, $s\le0$ forces $(a,s)=(2,0)$ and string length zero; in
dimension five it permits only $(a,s)=(2,-1)$, of length zero,
or $(2,0)$, of length one. Positive-index strings contribute no
kernel to the displayed maps. The two graded groups in the last
display and the grade-zero group in dimension four are not proved
to vanish here. This exact reduction sharpens the description of
the outstanding question; it is recorded as an equivalence, not as
a closure of the universal positivity problem.
